%% Generated by generate_slides.py. Edit this file or the source data, then run build.sh.
\begin{frame}{A short history of external memory}
\subt{Different traditions solve different parts of the memory problem.}
\begin{columns}[c,onlytextwidth]\column{.565\textwidth}
\begin{itemize}\setlength{\itemsep}{.42em}\footnotesize
\item \textbf{1945.} \chg{Vannevar Bush imagined the Memex: a desk that stores a scientist's books, papers and notes on microfilm and lets him save chains of links between them. It was never built.}
\item \textbf{2000s–2020s.} \chg{Personal knowledge management: real note-taking software with links between notes (Evernote, Roam, Obsidian), plus methods for filing notes, such as Tiago Forte's Projects, Areas, Resources, Archives (PARA). The person still does all the linking.}
\item \textbf{2010s–2020s.} Search, embeddings, retrieval-augmented generation and knowledge graphs.
\item \textbf{2024–2026.} \chg{The model stops only fetching documents and starts writing: an agent keeps an up-to-date summary of what the documents say (GraphRAG summaries, Karpathy's LLM Wiki).}
\medskip
\source{\chg{\href{https://www.theatlantic.com/magazine/archive/1945/07/as-we-may-think/3881/}{Bush (1945)}; \href{https://www.buildingasecondbrain.com/}{Forte (2022)}; \href{https://arxiv.org/abs/2005.11401}{Lewis et al. (2020)}; \href{https://arxiv.org/abs/2404.16130}{Edge et al. (2024)}; \href{https://gist.github.com/karpathy/442a6bf555914893e9891c11519de94f}{Karpathy (2026)}}}
\end{itemize}
\column{.405\textwidth}
\sboxblue{Avoid a false ladder}{RAG, wiki, graphs and notes coexist; no single approach makes the others obsolete.}
\sboxblue{Today’s question}{What persists across tasks: raw documents, distilled knowledge, decisions or executable worksflows?}
\end{columns}
\medskip
\tagline{A Brain architecture is a choice about what is stored, retrieved and updated.}
\end{frame}

\begin{frame}{Approach 1: put everything in the prompt}
\subt{The simplest external-memory strategy is often the least selective.}
\begin{columns}[c,onlytextwidth]\column{.565\textwidth}
\begin{itemize}\setlength{\itemsep}{.48em}\small
\item \textbf{Design.} Attach notes, instructions and all source material to each new request.
\item \textbf{Strength.} No retrieval stack, transparent inputs and straightforward debugging.
\item \textbf{Weakness.} Large inputs, repeated material and often low signal density.
\item \textbf{Best use.} Small and stable corpora for a narrow and high-stakes task where full context fits.
\end{itemize}
\column{.405\textwidth}
\sboxblue{\chg{Common mistake}}{\chg{Assuming that if the documents fit in the context window, the model reads all of them with equal care. Facts buried in the middle of a long input are recalled less often.}}
\sboxblue{Improvement}{Move from full read to indexed, evidence-backed selective disclosure.}
\end{columns}
\tagline{\chg{A larger window allow you to paste more, but it does not say which decision is current, which idea was rejected, or who approved a change.}}
\end{frame}

\begin{frame}{Approach 2: Retrieval-Augmented Generation (RAG}
\subt{RAG retrieves source evidence when a question arrives.}
\begin{itemize}\setlength{\itemsep}{.6em}
\item \textbf{Prepare.} Parse documents, split or structure text and build searchable indexes.
\item \textbf{Query.} Retrieve relevant passages using lexical search, vector similarity or both.
\item \textbf{Generate.} Answer with retrieved context and citations to original documents.
\item \textbf{Risks}:
\begin{itemize}
\item \chg{\textbf{Missed evidence.} Search never returns the appendix table, so the answer goes on without it.}
\item \chg{\textbf{Bad chunk boundaries.} A coefficient is cut off from the sentence that says which regression it belongs to.}
\item \chg{\textbf{Ambiguous provenance.} Two notes disagree, and the answer does not say which file it used.}
\item \chg{\textbf{Stale index.} The paper changed last week; search still returns last month's draft.}
\end{itemize}
\end{itemize}
\medskip
\tagline{\chg{RAG finds the passages in your documents that match a question. It does not keep track of where the project stands: current decisions, rejected ideas, next steps.}}
\end{frame}

\begin{frame}{Approach 3: Knowledge bases, links and GraphRAG}
\subt{\chg{An Obsidian link only says that two notes are related.\\
A knowledge graph also says how: this paper uses that dataset, this result contradicts that one.}}
\begin{columns}[c,onlytextwidth]\column{.565\textwidth}
\begin{itemize}\setlength{\itemsep}{.48em}\small
\item \textbf{Wiki and backlinks.} Human-readable notes form navigable relationships.
\item \textbf{Typed graph.} An edge can mean uses-method, produced-by or contradicts.
\item \textbf{GraphRAG.} Graph-based summaries can assist questions spanning many documents.
\item \textbf{Trade-off.} More expressive relationships require extraction, validation and maintenance.
\end{itemize}
\column{.405\textwidth}
\sboxblue{Economics graph}{Paper → identification → dataset → do-file → figure → reported effect.}
\sboxblue{Practical starting point}{First maintain linked Markdown; add typed relations only when necessary.}
\end{columns}
\end{frame}

\begin{frame}{Approach 3.1: A LLM-maintained wiki}
\subt{\chg{In April 2026 Andrej Karpathy (OpenAI, Anthropic) proposed that the agent write and keep up to date a wiki of what your sources say, so the next question starts from that summary instead of from the raw files.}}
\begin{itemize}\setlength{\itemsep}{.6em}
\item \textbf{Raw.} Immutable or preserved source materials: papers, web pages, transcripts.
\item \textbf{Wiki.} Interlinked, editable Markdown pages synthesizing what sources establish.
\item \textbf{Schema.} Agent-facing rules for naming, links, ingestion and maintenance.
\item \textbf{Compounding effect.} New sources update old pages; useful answers may be filed back into the wiki.
\end{itemize}
\medskip
\tagline{The agent is the maintainer; Obsidian is a readable interface.}
\end{frame}

\begin{frame}{LLM Wiki workflow: ingest, query, lint}
\subt{\chg{The cycle runs in this order: ingest new sources, query the wiki, lint it from time to time, promote what proved useful.}}
\begin{columns}[c,onlytextwidth]\column{.485\textwidth}
\sboxblue{\chg{1.} INGEST}{Preserve source; identify claims and entities; patch affected pages; preserve citations.}
\column{.485\textwidth}
\coursebox{\chg{2.} QUERY}{Read index; follow relevant pages; recover supporting sources; synthesize a reply.}
\end{columns}\vspace{1mm}
\begin{columns}[c,onlytextwidth]\column{.485\textwidth}
\sboxblue{\chg{3.} LINT}{Check broken links, orphan notes, duplications, contradictions and stale claims.}
\column{.485\textwidth}
\coursebox{\chg{4.} PROMOTE}{Save vetted synthesis or a reusable method for future tasks.}
\end{columns}\vspace{1mm}
\tagline{Record uncertain or contradictory evidence; do not silently replace it.}
\end{frame}

\begin{frame}{An effective hybrid: RAG + curated wiki}
\subt{Use one layer for detailed evidence and another for stable distilled knowledge.}
\begin{columns}[c,onlytextwidth]
\column{.47\textwidth}\begin{itemize}\setlength{\itemsep}{.5em}\small
\item \textbf{Original sources.} Keep papers, code, transcripts and records available for verification.
\item \textbf{Retrieval.} Search sparse, situational evidence when the task demands it.
\item \textbf{Curated memory.} Maintain dense, frequently used rules, positions, assumptions and project state.
\item \textbf{Governance.} Validate edits against source evidence before promoting them to canonical memory.
\end{itemize}
\column{.5\textwidth}
\begin{chgfig}
\scalebox{0.72}{\begin{tikzpicture}[x=1cm,y=1cm]
  \node[frameb,minimum width=5.6cm,minimum height=0.75cm,align=center] (s) at (0,2.6)
    {{\scriptsize\bfseries\color{ink}Original sources}\ {\scriptsize\color{muted}papers, data, code, transcripts}};
  \node[frameb,minimum width=2.6cm,minimum height=1.05cm,align=center] (r) at (-1.5,1.05)
    {{\scriptsize\bfseries\color{ink}Retrieval}\\[-1pt]{\tiny\color{muted}search when a task needs it}};
  \node[inkb,minimum width=2.6cm,minimum height=1.05cm,align=center] (w) at (1.5,1.05)
    {{\scriptsize\bfseries\color{ink}Curated memory}\\[-1pt]{\tiny\color{muted}state, decisions, wiki}};
  \node[accb,minimum width=2.6cm,minimum height=0.75cm,align=center,font=\scriptsize\bfseries,text=ink] (a) at (0,-0.55) {Agent answers};
  \node[frameb,minimum width=2.6cm,minimum height=0.75cm,align=center,font=\scriptsize\bfseries,text=ink] (h) at (0,-1.85) {You review};
  \draw[arr] (s.south -| r.north) -- (r.north);
  \draw[arr] (s.south -| w.north) -- (w.north);
  \draw[arr] (r.south) -- (a.north west);
  \draw[arr] (w.south) -- (a.north east);
  \draw[arr] (a.south) -- (h.north);
  \draw[arrA] (h.east) -- ++(1.65,0) |- node[pos=0.25,right,font=\tiny,text=accent,align=left]{promote\\what holds} (w.east);
\end{tikzpicture}}
\end{chgfig}
\end{columns}
\tagline{The retrieval layer complements memory; it does not replace it.}
\end{frame}

\begin{frame}{Which architecture should you choose?}
\subt{\chg{It depends on how much source material you have and how long the project lasts.}}
\begin{columns}[c,onlytextwidth]\column{.565\textwidth}
\begin{itemize}\setlength{\itemsep}{.48em}\small
\item \textbf{Prompt-only.} Small, finite task: minimal infrastructure, higher repetition cost.
\item \textbf{RAG.} Large reference corpus: stronger source discovery and evidence retrieval.
\item \textbf{LLM Wiki.} Repeated synthesis within a continuing domain: curated compounding knowledge.
\item \textbf{Hybrid.} Long-running research or enterprise work: state + wiki + targeted retrieval.
\end{itemize}
\column{.405\textwidth}
\sboxblue{\chg{Rule of thumb}}{\chg{A few dozen pages: paste them. Hundreds of documents: add search. A project that runs for months: keep files with its state and decisions.}}
\sboxblue{\chg{In this course}}{\chg{We build the hybrid: Markdown project files, plus search when the sources get large.}}
\end{columns}
\end{frame}
